\section{Conclusion}
\label{sec:conclusion}

This paper has presented a comprehensive analysis of NVIDIAScape (CVE-2025-23266), a critical container escape vulnerability arising from unsafe environment variable inheritance in the NVIDIA Container Toolkit's privileged OCI hook. We have examined the technical root cause (the absence of environment sanitisation at the \texttt{execve()} boundary) and traced the precise exploit chain by which an attacker within an unprivileged container can achieve arbitrary code execution with root privileges on the host.
 
The vulnerability exemplifies a class of security failures that recurs across the history of privileged system software: the unsafe consumption of untrusted input inherited through the process execution chain. In this specific case, the inherited input is the \texttt{LD\_PRELOAD} environment variable, and the dynamic linker's faithful processing of it becomes the mechanism of exploitation. The exploit requires no kernel vulnerability, no elevated container permissions, and no user interaction, only a standard GPU container and control over its environment variables.
 
The impact in modern AI infrastructure is severe. GPU nodes hosting multi-tenant LLM workloads represent some of the most data-rich and commercially sensitive targets in current cloud infrastructure, and the ubiquity of the NVIDIA Container Toolkit across GPU deployments means that the vulnerability had a very broad blast radius at the time of disclosure.
 
Several lessons emerge from this analysis. First, privileged system toolkits must treat their execution environment as untrusted when that environment is derived, directly or indirectly, from unprivileged input. Second, the Principle of Least Privilege must be applied not only to container workloads but equally to the privileged components of the container ecosystem itself. Third, service-based architectural designs for hardware management layers are inherently more robust than direct invocation models that inherit caller state. Fourth, defence-in-depth controls (Mandatory Access Control policies, seccomp profiles, rootless container runtimes) provide critical residual protection when individual components fail.
 
As AI workloads continue to demand GPU resources at scale, and as the infrastructure supporting them grows in complexity and strategic value, the security of the privileged software layer connecting containers to hardware will remain a critical and underexplored frontier. NVIDIAScape is a precise and instructive case study in both the fragility of this layer and the clear engineering principles that can make it more robust.